\documentclass[10pt,a4paper]{amsart}
\usepackage[margin=19mm]{geometry}
\usepackage{amsmath,amssymb,mathtools}
\usepackage[scr=rsfs,cal=euler]{mathalfa}
\usepackage{xfrac}
\usepackage[unicode,hidelinks,bookmarks=false]{hyperref}
\newcommand{\ie}{i.e.\ }
\newcommand{\cf}{cf.\ }
\newcommand{\resp}{respectively\ }
\newcommand{\Subsection}{\subsection}
\providecommand{\nobreakdash}{\nobreak\hbox{-}}
\setlength{\emergencystretch}{2em}
\multlinegap0pt
\usepackage{amssymb}
\usepackage{mathtools}
\usepackage{tikz}
\newtheorem{lemma}{Lemma}
\newcommand{\bE}{\mathbf{E}}
\newcommand{\e}{\mathbf{e}}
\def\rve{{\mathbf{e}}}
\def\rvz{{\mathbf{z}}}
\title{PAC-Bayesian Bounds for Learning Partially Observed Stochastic Linear Time-Invariant State-Space Systems with Inputs and Sub-Gaussian Noise}
\author{Mihaly Petreczky, Mohamad Al Ahdab, John Leth}
\date{Source: arXiv:2609.08740v1 (CC BY 4.0)}
\begin{document}
\maketitle

\noindent\emph{Source and licence.} PAC-Bayesian Bounds for Learning Partially Observed Stochastic Linear Time-Invariant State-Space Systems with Inputs and Sub-Gaussian Noise. Version arXiv:2609.08740v1 (CC BY 4.0), distributed by arXiv under the Creative Commons Attribution 4.0 International licence, http://creativecommons.org/licenses/by/4.0/, as recorded in the arXiv licence metadata for this exact version and shown on the abstract page https://arxiv.org/abs/2609.08740v1. Full licence notice: \texttt{licenses/arxiv-2609.08740-CC-BY-4.0.txt}.\par\smallskip

\noindent\textit{Editorial scope.} Counted unit: the article's lemma basic\_lemma:lipschitz (source0.tex lines 3428--3459). The article's complete original proof of it is delivered with it (source0.tex lines 3460--3564, one \texttt{proof} environment of 105 lines). No mathematics is written, repaired or re-derived: the statement and the proof are the article's own lines, in the article's own notation, and no retained line is substituted or re-flowed.  Retained uncounted as the source-local context the proof needs: lines 2992--3011; lines 3321--3332; lines 3394--3399.  Nothing was omitted from any retained span, so the package carries no omission record.  No retained line contains a citation, so the wrapper carries no bibliography and no bibliography entry.  No retained line contains a figure environment, an image-inclusion command or drawing code, so no image is delivered and the package declares no asset dependency.  Editorial formatting only, in addition: an AMS \texttt{amsart} preamble that restates the article's own declarations and macros used by the retained lines, namely the environments \texttt{lemma} on separate counters, together with the standard packages those lines need.  The retained environments print in this document as Lemma 1 in order of appearance, whereas the article prints the same items with the numbering of its own full document.  The complete native source of the article as distributed in the arXiv e-print for this exact version is bundled unchanged as \texttt{sources/209743/source0.tex}.
\begin{lemma}[Subgaussian noise]
\label{lemma:subgauss:Cramer}
If
$\|\e_g(t)\|_2$ is sub-Gaussian  then
%satisfies Cramer's condition, i.e.,
$\Psi_{\e_a}(\lambda)=\bE[e^{\lambda \|\e_g(0)\|_2}]$ 
%of the noise $\e_g(t)$
 is well-defined for all $\lambda > 0$.
%the moment generating function$\Psi_{\e}(\lambda)=\bE[e^{\lambda \|\e_g(0)\|_2}]$,
%is wel-defined for all $\lambda > 0$.
In fact, %there exists a constant $K_e \le C \sigma_e$, where $C$ is a universal
%constant such that
%with $C_{1,e}=\mu_{max}(Q_e)\sqrt{n_e!}$,
%and $C_{2,e}=\sqrt{n_e+2} \mu_{max}(Q_e)$
%$K_e=\max\{\mu_{max}(Q_e)^2n_e, C_{2,e}+C_{1,e}\}$
\begin{align*}
& \Psi_{\e_g}(\lambda) \le  \bar{\Psi}_{\rve_g}(\lambda)=e^{(\lambda^2 14^2 \|\rve_g\|_{\psi_2}^2+3\lambda \|\rve_g\|_{\psi_2})}
%1+\lambda C_{1,e} (1+\lambda C_{2,e})e^{\lambda^2 \mu_{max}(Q_e)^2 n_\e} \le  
\end{align*}
\end{lemma}

\begin{lemma}
\label{noise:boundC}
%With notation as in Lemma~\ref{lemma:subgauss:Cramer}
\begin{align*} 
     \bE[e^{s \|e_g(t)\|_2 \chi(\|\e_g(t)\|_2 > C)}]
     \le 
     \left(1+ s 
         \frac{\bar{\Psi}_{\rve_g}(\bar{\lambda})C}{e^{\bar{\lambda}C}-1} + s^2
         \frac{2\bar{\Psi}_{\rve_g}(\bar{\lambda})}{\bar{\lambda}^2}\right):=\Psi_{\rve_g}(s,\bar{\lambda},C)
 \end{align*}
 for all $s \in [0,0.5\bar{\lambda})$. 
\end{lemma}

\begin{lemma}
\label{basic_lemma:lipschitz:lemma_aux}
  Let $a_k \ge 0$ be a sequence of positive numbers such that the series $\sum_{k=0}^{\infty} a_k$ is convergent. For any $N$, define
  $c_{k,N}=\sum_{l=\max\{0,k-N+1\}}^k a_l$. Then  $\sum_{k=0}^{L} c_{k,N} \le N \sum_{k=0}^{L} a_k$, 
  the series $\sum_{k=0}^{\infty} c_{k,N}$ is convergent, and $\sum_{k=0}^{\infty} c_{k,N} \le N \sum_{k=0}^{\infty} a_k$. 
\end{lemma}

\begin{lemma}
\label{basic_lemma:lipschitz}
% Let $\mathbf{I}$ be an index set, and let $\{\rvz_i\}_{i \in \mathbf{I}}$ be a collection of stochastic process
% such that for all $i \in \mathbb{I}$, $\rvz_{i}(t)=(\beta_{i,t} \star \e_g)(t)$ and
% $\bar{\rvz}(t)=(\beta_{i,t} \star \bar{\e}_g)(t)$
% %\sum_{k=0}^{\infty} \beta_{i,t}(k) \bar{\e}_g(t-k)$, for 
% suitable matrices $\{\beta_{i,t}(k)\}_{k=0}^{\infty}$, $i \in \mathbf{I}$, 
% such that %$\sup_{i \in \mathbf{I}} \|\beta_i\|_{\ell_1} <  +\infty$, where
% %for 
% $\|\bar{\beta}\|_{\ell_1}< +\infty$, where $\bar{\beta}(k)=\sup_{i \in \mathbf{I},t \in \mathbb{Z}} \|\beta_{i,t}(k)\|_2$.
% %$\bar{H}(\underline{\beta}):=\sum_{k=0}^{\infty}   +\infty$. 
Let $\mathbf{I}$ be an index set, and $\{\beta_{i,t}(k)\}_{k=0}^{\infty}$, $i \in \mathbf{I}$, be a sequence of matrices such that $\|\bar{\beta}\|_{\ell_1}< \infty$, where $\bar{\beta}(k)=\sup_{i \in \mathbf{I},t \in \mathbb{Z}} \|\beta_{i,t}(k)\|_2$. Let $\rvz_{i}(t)=(\beta_{i,t} \star \e_g)(t)$ and
$\bar{\rvz}_i(t)=(\beta_{i,t} \star \bar{\e}_g)(t)$.
Then for $\lambda < 0.5N$
 \begin{align*}
 & \bE[e^{\frac{\lambda}{N} \sup_{i \in \mathbf{I}}\sum_{t=0}^{N-1} \|\rvz(t)-\bar{\rvz}_i(t)\|_2}]
  \le \prod_{k=0}^{\infty} 
  \bar{\Psi}_{\rve_g}\left(\frac{c(k) \lambda}{N},\|\bar{\beta}\|_{\ell_1},C\right) \le 
   e^{\mathcal{C}(\|\bar{\beta}\|_{\ell_1},\lambda,N,C)}  \\
   & \mathcal{C}(s,\lambda,N,C)=\lambda \frac{\bar{\Psi}_{\rve_g}(s)s^2C}{(e^{sC}-1)}+
     \frac{\lambda^2}{N} 2\bar{\Psi}_{\rve_g}( s)
  %(1+\frac{\lambda c_{k,g}}{N} 
  %    \frac{\Phi(\|\Sigma_{gen}\|C) \|\Sigma_{gen}\|}{e^{\|\Sigma_{gen}\|C}-1} + \frac{\lambda^2 c_{k,g}^2}{N^2}
  %    \frac{2\Phi(\|\Sigma_{gen}\|C)}{\|\Sigma_{gen}\|^2})
 \end{align*} 
 where 
 %$c_{k}=\sup_{i \in \mathbf{I}} c_k(i)$ and  
 $c(k)=\sum_{l=\max\{0,k-N+1\}}^{k} \|\bar{\beta}(l)\|_2$
 and the infinite product on the right-hand side is well-defined. 
 %$\alpha_{k,g}=C_gA_g^{k-1}K_g$, $k > 1$
%$\alpha_{0,g}=I_{n_y}$.
\end{lemma}

\begin{proof}[Proof of Lemma \ref{basic_lemma:lipschitz}]
  Note that 
  \begin{align}
    %& c_{i,t}(k) \le \sum_{l=0}^{\infty} \|\beta_{i,t}(l)\|_2 \le \|\bar{\beta}\|_{\ell_1} \label{basic_lemma:lipschitz:eq0.1} \\
    &  c(k) 
    %= \sum_{l=\max\{0,k-N+1\}}^{k} \bar{\beta}(l) 
    \le \|\bar{\beta}\|_{\ell_1}
      \label{basic_lemma:lipschitz:eq0.2} 
   \end{align}
    and by Lemma \ref{basic_lemma:lipschitz:lemma_aux}
    \begin{equation}
      \label{basic_lemma:lipschitz:eq1}
      \begin{aligned}
     & \sum_{k=0}^{\infty} c(k) \le N \sum_{k=0}^{\infty} \|\bar{\beta}(k)\|_2 = N \|\bar{\beta}\|_{\ell_1} \\
     %& \sum_{k=0}^{\infty} \sup_{i \in \mathbf{I}} c_{i}(k) \le N \sum_{k=0}^{\infty} \sup_{i \in \mathbf{I}}  \|\beta_i(k)\|_2 = N \|\bar{\beta}\|_{\ell_1} \\
      \end{aligned}
    \end{equation}
    %Indeed, the first inequality is a direct consequence of Lemma \ref{basic_lemma:lipschitz:lemma_aux}, while the second one follows
    %from the statement of Lemma \ref{basic_lemma:lipschitz:lemma_aux} 
    %and \eqref{basic_lemma:lipschitz:eq0.2}.
    Note now that 
      \begin{align*}
      & \sum_{t=0}^{N-1}  \|\rvz_{i,t}(t)-\bar{\rvz}_{i,t}(t)\|_2  \\
      & \le
      \sum_{t=0}^{N-1}  \sum_{k=0}^{\infty} \|\beta_{i,t}(k)\|_2 \|\e_g(t-k)-\bar{\e}_g(t-k)\|_2 \\
      & \le \sum_{t=0}^{N-1}  \sum_{k=0}^{\infty} \|\bar{\beta}(k)\|_2 \|\e_g(t-k)\chi(\|\e_g(t-k)\|_2 > C)\|_2 \\
      & =\sum_{k=0}^{\infty} c(k) \|\e_g(N-k-1) \chi(\|\e_g(N-k-1)\|_2 > C)\|_2 
        \end{align*}    
        Hence,
        \begin{align*}
      &\bE[ e^{\frac{\lambda}{N} \sup_{i \in \mathbf{I}, t \in \mathbb{Z}} \sum_{t=1}^{N} \| \|\rvz_{i,t}(t)-\bar{\rvz}_{i,t}(t)\|_2}]
      \le \\
      %& \bE[e^{\frac{\lambda}{N} \sup_{i \in \mathbf{I}, t \in \mathbb{Z}}\sum_{k=0}^{\infty} c_{i,t}(k) \|e_g(N-k-1) \chi(\|\e_g(N-k-1)\|_2 > C)\|_2}] \\
      & \le \bE[e^{\frac{\lambda}{N} \sum_{k=0}^{\infty} 
      %\sup_{i \in \mathbf{I}, t \in \mathbb{Z}} c_{i,t}(k) 
      c(k) \|\e_g(N-k-1) \chi(\|\e_g(N-k-1)\|_2 > C)\|_2}] \\
      & = \lim_{L \rightarrow \infty} 
      \prod_{k=0}^{L} \bar{\Psi}_{\rve_g}\left(\frac{c(k) \lambda}{N},\|\bar{\beta}\|_{\ell_1},C\right)
        \end{align*}
        where we applied  Lemma \ref{noise:boundC}
        to $\bE[e^{\frac{\lambda}{N} c(k)\|e_g(N-k-1) \chi(\|\e_g(N-k-1)\|_2 > C)\|_2}]$
        with $\bar{\lambda}=\|\bar{\beta}\|_{\ell_1}$ and $s=\frac{c(k)\lambda}{N} < 0.5\|\bar{\beta}\|_{\ell_1}$ since $\lambda<0.5N$ by assumption.

        Note that 
        \[
        \begin{aligned}
        &  d_L=\ln\left(\prod_{k=0}^{L} \bar{\Psi}_{\rve_g}\left(\frac{c(k) \lambda}{N},\|\bar{\beta}\|_{\ell_1},C\right) \right)=\sum_{k=0}^{L} r_k \\
        & r_k:=\ln \left(1+\frac{c(k) \lambda}{N} \frac{\bar{\Psi}_{\rve_g}(\|\bar{\beta}\|_{\ell_1})\|\bar{\beta}\|_{\ell_1}C}{e^{\|\bar{\beta}\|_{\ell_1}C}-1} + \left(\frac{c(k) \lambda}{N}\right)^2 \frac{2\bar{\Psi}_{\rve_g}(\|\bar{\beta}\|_{\ell_1})}{\|\bar{\beta}\|_{\ell_1}^2} \right) \geq 0
        \end{aligned}
        \]
        %As $c(k) \ge 0$, $\|\bar{\beta}\|_{\ell_1} > 0$, $r_k \ge 0$. 
        We show that the series $\sum_{k=0}^{\infty} r_k$ 
        is bounded, more precisely, 
        \[
        \sum_{k=0}^{L} r_k \le 
      \|\bar{\beta}\|_{\ell_1} \lambda \frac{\bar{\Psi}_{\rve_g}(\|\bar{\beta}\|_{\ell_1})\|\bar{\beta}\|_{\ell_1}C}{e^{\|\bar{\beta}\|_{\ell_1}C}-1} + \frac{\lambda^2}{N} (\|\bar{\beta}\|_{\ell_1})^2 \frac{2\bar{\Psi}_{\rve_g}(\|\bar{\beta}\|_{\ell_1})}{\|\bar{\beta}\|_{\ell_1}^2} = \mathcal{C} (\|\bar{\beta}\|_{\ell_1},\lambda,N,C)
        \]
        From this it follows that $d_L$ is convergent, and from
        \[
      \prod_{k=0}^{ L} \bar{\Psi}_{\rve_g}\left(\frac{c(k) \lambda}{N},\|\bar{\beta}\|_{\ell_1},C\right) = e^{d_L}
        \]
        it follows that  
        \begin{equation}
       \label{basic_lemma:lipschitz:eq1.3}
       \begin{aligned}
        & \lim_{L \rightarrow \infty} 
        \prod_{k=0}^{ L} \bar{\Psi}_{\rve_g}\left(\frac{c(k) \lambda}{N},\|\bar{\beta}\|_{\ell_1},C\right) \le  e^{\mathcal{C}(\|\bar{\beta}\|_{\ell_1},\lambda,N,C)}
       \end{aligned}
        \end{equation}
        i.e., the infinite product is well-defined, and the statement of the lemma holds. 

        In order to show \eqref{basic_lemma:lipschitz:eq1.3}, note that using $\ln(1+x) \le x$, it follows that 
        \[
        \begin{aligned}
      & r_k \le \frac{c(k) \lambda}{N} \frac{\bar{\Psi}_{\rve_g}(\|\bar{\beta}\|_{\ell_1})\|\bar{\beta}\|_{\ell_1}C}{e^{\|\bar{\beta}\|_{\ell_1}C}-1} + \left(\frac{c(k) \lambda}{N}\right)^2 \frac{2\bar{\Psi}_{\rve_g}(\|\bar{\beta}\|_{\ell_1})}{\|\bar{\beta}\|_{\ell_1}^2}
        \end{aligned}
        \]
        Hence,
        \[
        \begin{aligned}
      & \sum_{k=0}^{L} r_k  \le \frac{\lambda \sum_{k=0}^{L}c(k)}{N} \frac{\bar{\Psi}_{\rve_g}(\|\bar{\beta}\|_{\ell_1})\|\bar{\beta}\|_{\ell_1}C}{e^{\|\bar{\beta}\|_{\ell_1}C}-1} +  \left(\frac{\sum_{k=0}^{L} (c(k))^2}{N^2}\right) \frac{2\bar{\Psi}_{\rve_g}(\|\bar{\beta}\|_{\ell_1})}{\|\bar{\beta}\|_{\ell_1}^2}
        \end{aligned}
        \]
        
        From \eqref{basic_lemma:lipschitz:eq0.2}
        it follows that
        $c(k) \le \|\bar{\beta}\|_{\ell_1}$.
        %Moreover, 
        %\[ (\frac{c_{i,t}(k)}{N})^2=\left(\frac{1}{N} \sum_{l=\max\{0,k-N+1\}}^k \beta_{i,t}(l) \right)^2 \le \frac{1}{N} \sum_{l=\max\{0,k-N+1\}}^k \beta_{i,t}(l)^2 \le \frac{1}{N} \sum_{l=\max\{0,k-N+1\}}^k (\bar{\beta}(l))^2
        %\]
        Hence
        \[
         \left(\frac{1}{N} c(k)\right)^2 \le \frac{1}{N^2} c(k)  \|\bar{\beta}\|_{\ell_1}
        \]
        and as from \eqref{basic_lemma:lipschitz:eq1} it follows that $\sum_{k=0}^{L} c(k) \le N \|\bar{\beta}\|_{\ell_1}$, we get that 
        % as by Lemma \ref{basic_lemma:lipschitz:lemma_aux}
        %
        \begin{align*}
        & \sum_{k=0}^{L} \left(\frac{1}{N} c(k)\right)^2 \le 
        \frac{1}{N^2} \sum_{k=0}^{L}  c(k)
        \|\bar{\beta}\|_{\ell_1} \le 
        \frac{1}{N} \|\bar{\beta}\|_{\ell_1}^2
      \end{align*}
      This completes the proof.
    \end{proof}

\end{document}
